\documentclass[border=8pt]{standalone}
\usepackage[utf8]{inputenc}
\usepackage[T1]{fontenc}
\usepackage{lmodern}
\usepackage[spanish,es-nodecimaldot]{babel}
\usepackage{siunitx}
\sisetup{output-decimal-marker={,}}
\usepackage[european,straightvoltages]{circuitikz}
\usetikzlibrary{babel,calc}
\definecolor{accent}{HTML}{2E5C8A}
\definecolor{graytext}{HTML}{595959}
\begin{document}
\begin{circuitikz}[font=\scriptsize]
  % ================= Etapa 1 (TL082 A: patas 1, 2, 3) =================
  \draw (0,0) node[left]{$V_s$ (divisor LDR)} node[ocirc]{}
        to[eC, l={$C_1$ \qty{10}{\micro\farad}}] (2.4,0) coordinate(a);
  \draw (a) to[R, l_={$R_{B1}$ \qty{47}{\kilo\ohm}}, *-] (2.4,2.1) node[above]{GND};
  \node[op amp, noinv input up, anchor=+] (U1) at (3.6,0) {};
  \draw (a) -- (U1.+);
  \draw (U1.-) -- ++(-0.4,0) -- ++(0,-1.4) coordinate(p);
  \draw (p) to[R, l_={$R_1$ \qty{4.7}{\kilo\ohm}}] ++(-2.2,0) node[left]{GND};
  \draw (p) to[R, l={$R_2$ \qty{100}{\kilo\ohm}}, *-] (p -| U1.out) coordinate(r2) -- (U1.out);
  \draw (p) -- ++(0,-1.1) coordinate(q) to[C, l_={$C_f$ \qty{150}{\nano\farad}}] (q -| U1.out) -- (r2);
  \node[font=\tiny, graytext] at ($(U1.+)+(0.2,0.22)$) {3};
  \node[font=\tiny, graytext] at ($(U1.-)+(0.2,-0.22)$) {2};
  \node[font=\tiny, graytext] at ($(U1.out)+(0.05,-0.25)$) {1};
  \node[graytext] at ($(U1.center)+(0,0.9)$) {TL082 A};
  % ================= Etapa 2 (TL082 B: patas 5, 6, 7) =================
  \draw (U1.out) to[short,*-] ++(0.3,0) node[above]{$v_1$}
        to[eC, l={$C_2$ \qty{10}{\micro\farad}}] ++(2.2,0) coordinate(b);
  \draw (b) to[R, l_={$R_{B2}$ \qty{47}{\kilo\ohm}}, *-] (b |- 0,2.1) node[above]{GND};
  \node[op amp, noinv input up, anchor=+] (U2) at ($(b)+(1.2,0)$) {};
  \draw (b) -- (U2.+);
  \draw (U2.-) -- ++(-0.4,0) -- ++(0,-1.4) coordinate(p2);
  \draw (p2) to[R, l_={$R_3$ \qty{10}{\kilo\ohm}}] ++(-2.2,0) node[left]{GND};
  \draw (p2) to[vR, l_={$R_4$ trimmer \qty{100}{\kilo\ohm}}, *-] (p2 -| U2.out) -- (U2.out);
  \node[font=\tiny, graytext] at ($(U2.+)+(0.2,0.22)$) {5};
  \node[font=\tiny, graytext] at ($(U2.-)+(0.2,-0.22)$) {6};
  \node[font=\tiny, graytext] at ($(U2.out)+(0.05,-0.25)$) {7};
  \node[graytext] at ($(U2.center)+(0,0.9)$) {TL082 B};
  % ================= Desplazamiento a 2,5 V y protección de A0 =================
  \draw (U2.out) to[short,*-] ++(0.3,0) node[above]{$v_2$}
        to[eC, l={$C_3$ \qty{10}{\micro\farad}}, invert] ++(2.2,0) coordinate(c);
  \draw (c) to[R, l_={\qty{100}{\kilo\ohm}}, *-] ++(0,2.1) node[above]{5V (Nano)};
  \draw (c) to[R, l={\qty{100}{\kilo\ohm}}] ++(0,-2.1) node[ground]{};
  \draw (c) to[R, l={$R_o$ \qty{10}{\kilo\ohm}}] ++(2.4,0) coordinate(d);
  \draw (d) to[D, l_={$D_1$}, *-] ++(0,2.1) node[above]{5V (Nano)};
  \draw (d) to[D, l={$D_2$}, invert] ++(0,-2.1) node[ground]{};
  \draw (d) -- ++(1.8,0) coordinate(e) to[short,*-o] ++(0.9,0) node[right]{\textbf{A0}};
  \draw (e) to[C, l={$C_o$ \qty{100}{\nano\farad}}] ++(0,-2.1) node[ground]{};
  % ================= Alimentación =================
  \node[draw=accent, dashed, rounded corners, align=left, anchor=north west, font=\scriptsize, inner sep=5pt]
    at (0.2,-4.4) {%
    \textbf{Alimentación TL082:} pata 8 $\rightarrow$ $+\qty{12}{V}$, pata 4 $\rightarrow$ $-\qty{12}{V}$
    (fuente doble del laboratorio en serie).\\
    \qty{100}{\nano\farad} de pata 8 a GND y de pata 4 a GND, junto al chip.
    \textbf{COM de la fuente = GND del Nano.}\\
    $D_1$, $D_2$: 1N4148 (limitan A0 a $-0{,}6\ldots5{,}6$ V). $C_1$, $C_2$, $C_3$ electrolíticos: el $+$ hacia la izquierda en $C_1$; hacia la derecha en $C_3$.};
\end{circuitikz}
\end{document}
